\id{ҒТАМР 50.41.25}{}

{\bfseries RUBERT-ДИАГНОСТИКАСЫ ЖӘНЕ IRT-АДАПТИВТІ ТЕСТІЛЕУ НЕГІЗІНДЕ SQL
ЖЕКЕЛЕНДІРІЛГЕН ОҚЫТУ}

{\bfseries \tsp{1}Б.Ж.
Жарлықасов}\fig{i/image1}[]{\bfseries ,
\tsp{2}Г.М.
Баенова}\fig{i/image1}[]{\bfseries ,
\tsp{1}Қ.С.
Мауленов}\fig{i/image1}[]{\bfseries ,
\tsp{2}А.М.
Сыздыков}\fig{i/image1}[]\envelope 

\emph{\tsp{1}Ахмет Байтұрсынұлы атындағы Қостанай өңірлік
университеті, Қостанай, Қазақстан,}

\emph{\tsp{2}Л.Н. Гумилев атындағы Еуразия ұлттық
университеті, Астана, Қазақстан}

\envelope Корреспондент-автор:
aierke91@mail.ru

Білім беру технологиялары мен оқу процесінде цифрлық платформалар және
орталар ажырамас бөлікке айналуда. Оқу материалдары мен тапсырмалар
білім алушыларға цифрлық форматта дайындалып ұсынылғандықтан, ол білім
алу және оқыту процесіне елеулі өзгерістер енгізуде. Бұл өзгерістер тек
цифрлық форматтағы оқыту әдістеріне ғана емес, сонымен қатар білімді
бағалауға да қатысты.

Осы жұмыста SQL тілін қолдауға арналған кешенді жасанды интеллект жүйесі
ұсынылған. Ол үш негізгі компонентті біріктіреді: RUBERT моделіне
негізделген студенттердің қателерін түрлері бойынша жіктеу модулі,
XGBoost алгоритіміне негізделген ұсынымдық модуль бұл студенттердің
қателер профилі мен жұмыс қарқынына сүйене отырып қайталау тақырыбын
анықтайды, тапсырмаларға жауап беру теориясының (IRT) екі параметрлі
моделін қолданатын бейімделгіш тестілеу жүйесі, ол білім алушының
ағымдағы қабілет деңгейіне сәйкес сұрақтардың күрделілігін динамикалық
түрде таңдайды. Жүйе веб-қосымша түрінде жүзеге асырылған және
«Ақпараттық жүйелер» бағытындағы 120 студенттің: SQL логтары, теориялық
сұрақтарға жауаптары және тест тапсырудың уақыт белгілері сияқты
деректер негізінде оқу қосымшасы (2024 жылдың қыркүйегі мен 2025жылдың
мамыры аралығында) қолдану барысында сыналды.

{\bfseries Түйін сөздер}: цифрлық оқыту платформалары, жасанды интеллект,
машиналық оқыту, оқыту модельдері, бағалау көрсеткіштері.

{\bfseries ПЕРСОНАЛИЗИРОВАННОЕ ОБУЧЕНИЕ SQL НА ОСНОВЕ RUBERT-ДИАГНОСТИКИ И
IRT-АДАПТИВНОГО ТЕСТИРОВАНИЯ}

{\bfseries \tsp{1}Б.Ж. Жарлыкасов, \tsp{2}Г.М.
Баенова, \tsp{1}К.С. Мауленов, \tsp{2}А.М.
Сыздыкова\envelope }

\emph{\tsp{1}Костанайский региональный университет им. Ахмет
Байтұрсынұлы, Костанай, Казахстан,}

\emph{\tsp{2}Евразийский национальный университет им. Л.Н.
Гумилева, Астана, Казахстан,}

e-mail: aierke91@mail.ru

В образовательных технологиях и учебном процессе цифровые платформы и
среды стали неотъемлемой частью. Учебные материалы и задания оформляются
и предоставляются обучающимся в цифровом формате, что внесло
существенные изменения в процесс обучения и освоения знаний. Эти
изменения касаются не только методов обучения в цифровом формате, но и
проверки знаний.

В работе представлена комплексная ИИ-система поддержки обучения SQL,
которая объединяет три ключевых компонента: классификация студенческих
ошибок по типам, основанная на модели RuBERT, рекомендательный модуль на
XGBoost для выбора тем повторения на основе профиля ошибок и темпа
работы, и адаптивное тестирование с применением двухпараметрической
модели теории ответов на задания (IRT), позволяющее динамически
подбирать сложность вопросов под текущую оценку способности обучаемого.
Система реализована в виде веб-приложения и протестирована на данных 120
студентов направления «Информационные системы», собранных в ходе
естественного использования учебного приложения (сентябрь 2024 -- май
2025): SQL-логи, ответы на теоретические вопросы и временные метки
прохождения тестов.

{\bfseries Ключевые слова}: цифровые обучающие платформы, искусственный
интеллект, машинное обучение, модели обучения, метрики оценки.

{\bfseries PERSONALIZED SQL TRAINING BASED ON RUBERT DIAGNOSTICS AND
IRT-ADAPTIVE TESTING}

{\bfseries \tsp{1}B.Zh. Zharlykassov, \tsp{2}G.M.
Baenova, \tsp{1}K.S. Maulenov, \tsp{2}A.M.
Syzdykova\envelope }

\emph{{\bfseries \tsp{1}}Akhmet Baitursynuly Kostanay Regional
University, Kostanay, Kazakhstan,}

\emph{{\bfseries \tsp{2}}L.N. Gumilyov Eurasian National
University, Astana, Kazakhstan,}

e-mail: aierke91@mail.ru

Digital platforms and environments have become an integral part of
educational technologies and the learning process. Learning materials
and assignments are designed and provided to students in digital format,
which has made significant changes to the process of learning and
acquiring knowledge. These changes concern not only digital teaching
methods but also knowledge testing.

The paper presents a comprehensive AI-based SQL training support system
that combines three key components: a classification of student errors
by type based on the RuBERT model, an XGBoost-based recommendation
module for selecting revision topics based on error profile and work
pace, and adaptive testing using a two-parameter item response theory
model (IRT), allowing dynamic adjustment of the complexity of questions
to the current assessment of the student' s ability. The
system is implemented as a web application and tested on data from 120
students majoring in Information Systems, collected during their natural
use of the learning application (September 2024 - May 2025): SQL logs,
answers to theoretical questions, and timestamps of passing tests.

{\bfseries Keywords:} digital learning platforms, artificial intelligence,
machine learning, learning models, assessment metrics.

{\bfseries Кіріспе.} Жоғары білім берудің цифрландырылуы нәтижесінде
электрондық білім беру ортасы мен оқытуды басқару жүйелері (LMS)
университеттердің негізгі инфрақұрылымына айналды. Оқу материалдары,
тапсырмалар мен бақылау іс-шаралары бүгінде цифрлық форматта дайындалады
және өткізіледі, ал оқу қызметі туралы жинақталған деректер оқытудың
моделін дәлірек құруға мүмкіндік береді. Сонымен қатар, күнделікті
операцияларды автоматтандыруға және жекелендіру мүмкіндіктерін кеңейтуге
қабілетті жасанды интеллект әдістері қарқынды дамып келеді.
Интеллектуалды платформаның негізі ретінде «Деректер базасы» пәні
таңдалды. Себебі біріншіден, бұл пән - ІТ мамандарын даярлаудың және
пәнаралық бағдарламалардың (деректер аналитикасы, ақпараттық жүйелер,
киберқауіпсіздік, бағдарламалық қамтамасыз ету инженериясы) іргелі
құрамдас бөлігі болып табылады, екіншіден, мақала авторлары осы пәнді өз
университеттерінде оқытады.

Транзакциялық жүйелер мен электрондық коммерциядан бастап, аналитикалық
витриналар мен шешім қабылдауды қолдау жүйелеріне дейін деректер
базалары қазіргі цифрлық экономиканың негізін құрайды. Негізгі
құзыреттерге реляциялық модельді түсіну, схемаларды жобалау,
нормализациялау, SQL тілін меңгеру, транзакциялар мен блокировкалар
механизмдері, индекстеу және сұраныстарды оңтайландырудың негізгі
қағидаттары жатады. Аталған тақырыптарды игеру сапасы түлектердің
қолданбалы міндеттерді шешу қабілетіне тікелей әсер етеді, ал
университеттер үшін білім беру бағдарламаларының өзектілігін және
студенттердің академиялық табыстылығын қамтамасыз ету тұрғысынан
маңызды.

«Деректер базасы» пәні магистратураға қабылдау емтихандарында ерекше
орын алады. Емтихандар мен біліктілік тесттері тек теориялық білімді
ғана емес, сонымен қатар практикалық дағдыларды да тексереді: яғни дұрыс
SQL-сұраныстарын құра білу, тапсырма шектеулеріне сәйкес деректер
құрылымын жобалау, орындау жоспарларын талдау және қателердің көздерін
анықтау қабілеттерін бағалайды. Алайда дәстүрлі бақылау түрлері (жауап
нұсқасын таңдау, терминдерді білуге арналған тапсырмалар) жобалау және
аналитикалық ойлау дағдыларының қалыптасу деңгейін толық көрсете
алмайды. Сондықтан әдістемелік міндет туындайды: теориялық білімді
тексеруді практикалық дағдыларды диагностикалаумен ұштастыру, сондай-ақ
студентке тақырыптар бойынша олқылықтарды көрсететін мазмұнды кері
байланыс ұсыну. Дегенмен, білім беру саласында жасанды интеллектке деген
қызығушылықтың артуына қарамастан, деректер базаларын оқыту саласы әлі
де жеткілікті зерттелмеген. Әсіресе, адаптивті оқыту мен студенттердің
қателіктерін талдау контексі бағыттында зерттеулер аз {[}1, 2{]}. Бұл
жағдай нақты білім беру міндеттерін ескере отырып, SQL оқыту процесіне
жасанды интеллект технологияларын интеграциялауға бағытталған қолданбалы
шешімдерге деген қажеттілікті арттырады.

Іс жүзінде LMS (мысалы, Moodle) және жаппай ашық білім беру
платформалары (мысалы, OpenEdu) контентті басқару, кесте құру, негізгі
есеп беру және қолжетімділік мәселелерін тиімді шешеді. Дегенмен,
«Деректер базасы» пәні үшін мұндай құралдар бірқатар себептерге
байланысты жеткіліксіз болып табылады.

Біріншіден, стандартты конфигурациядағы тестілер негізінен дұрыс жауапты
тануға бағытталған, яғни шешімді қайта жасаудың және құрудың орнына,
дұрыс жауапты тану қабілетін (яғни «тану» түріндегі тексеруді)
бағалайды. Бұл үстірт оқыту стратегиясын және іздеу әрекетін (ұсынылған
материалдардан жауапты тез табу дағдысын) қалыптастырады, бірақ
сұраныстарды дұрыс құра білу және олардың дұрыстығын түсіндіру қабілетін
дұрыс бағаламайды. SQL-ді тексеруге арналған плагиндер болса да, оларды
енгізу жеке баптауды, тапсырмалар мен тестілік базалардың сапалы
жиынтығын, сондай-ақ нақты тәжірибеде әрдайым жүзеге асырыла бермейтін
әдістемелік қолдауды қажет етеді.

Екіншіден, қолданыстағы тестілеу модульдері көбіне ішкі тақырыптар
деңгейінде диагностикалық талдау ұсынбайды: оқытушыларға топтың нақты
қай тұста қиналып жатқанын (мысалы: JOIN, агрегаттау, ішкі сұраулар,
терезелік функциялар) тез анықтау қиын, ал студент қателердің себептері
жөнінде егжей-тегжейлі түсіндірмелер мен жеке ұсыныстарға қол жеткізе
алмайды. Сонымен бірге, кері байланыс көбінесе екілік сипатта болады
(дұрыс/бұрыс) және кезең-кезеңмен түсіндіру механизмін қамтымайды. Бұл
бақылаудың педагогикалық құндылығын төмендетіп, қателерді түзету циклін
ұзартады.

Үшіншіден, шектеулі адаптивтілік: оқыту траекториясы белгілі бір
студенттің қателіктерінің профиліне нашар бейімделеді, білімнің қазіргі
жағдайына байланысты тапсырмалардың күрделілігі мен форматын автоматты
түрде реттеу мүмкіндігі жоқ. Жаппай платформалар (соның ішінде OpenEdu
курстары да) көбінесе жеке диагностика мен олқылықтарды кезең-кезеңімен
түзетуге емес, мазмұн мен қорытынды бағалауға бағытталған. Нәтижесінде
цифрлық ортаның мүмкіндіктері мен пәннің нақты қажеттіліктері арасында
алшақтық пайда болады, мұнда дұрыс шешімдерді құру қабілетін сыни
тұрғыдан бағалау қажет.

SQL-сұраныстарды автоматтандырылған тексеру бойынша жүргізілген
зерттеулерде сұраныстар сапасын бағалаудың әртүрлі тәсілдерін
қарастырады. Бір тәсіл SQL‑сұраныстарды ішінара бағалауды (partial SQL
query assessment) зерттейді, бұл тек "дұрыс/қате" деп қана анықтауға
емес, сонымен қатар эталондық нұсқамен құрылымдық ұқсастық деңгейін
бағалауға мүмкіндік береді. Тағы бір тәсіл SQL-ді автоматты түрде
бағалау үшін назар аудару механизмі бар конволюциялық нейрондық
желілерді пайдалануға негізделген {[}3, 4{]}. Бұл зерттеулер осы жұмыста
ұсынылған жүйе сәйкес келетін тақырыптың өзектілігі мен белсенді дамуын
көрсетеді.

Интернет дәуірі мен тілдік модельдердің кеңінен қолжетімділігі жаңа
қиындықтарды алдыңғы қатарға шығарды. Білім берудегі жасанды интеллект
оқыту үдерісін түбегейлі түрлендіре алады: \{ол білім алушыларға
жекелендірілген оқу материалдарын, автоматтандырылған кері байланысты
және адаптивті бағалау жүйелерін ұсына отырып, оқу стратегияларын, оқыту
сапасын және жалпы білім нәтижелерін жақсартуға мүмкіндік береді
{[}5{]}. Жасанды интеллект өзінің айтарлықтай әлеуетіне қарамастан,
когнитивтік парадоксқа да ие: ол оқу процесін жетілдіріп қана қоймай,
сонымен қатар білім алушылардың когнитивтік қабілеттерін, әсіресе сыни
тұрғыдан ойлау мен мәселелерді шешу дағдыларын әлсіретуі де мүмкін
{[}6{]}. Сондықтан жасанды интеллект технологияларын енгізу сақтықты,
әрі ойластырылған әдіс-тәсілдерді талап етеді, сонымен қатар басты
назарды тек рутиналық процестерді автоматтандыруға емес, жоғары деңгейлі
когнитивтік дағдыларды дамытуға бағыттаған дұрыс. Технологияға шамадан
тыс тәуелділікті болдырмау үшін автоматтандырылған қолдау мен тәуелсіз
танымдық қабілеттерді дамыту арасындағы тепе-теңдікті сақтау маңызды
болып табылады {[}7{]}.

Бір жағынан, студенттер код үлгілеріне және дайын SQL фрагменттеріне
жылдам қол жеткізу мүмкіндігіне ие, бұл олардың оқу үдерісін едәуір
жеделдетуге мүмкіндік береді. Екінші жағынан, студенттердің өздігінен
ойлау қабілетінің механикалық көшірумен алмастырылу қаупі артып,
деректер схемасын талдау, тұтастық шектеулерін түсіну және нәтиженің
дұрыстығын бағалау талап етілетін тапсырмаларға төзімділіктің
төмендеуіне алып келеді. Оқытушылар үшін бұл формальды тұрғыдан дұрыс,
бірақ тұжырымдамалық жағынан қате шешімдерді ажырату қажеттілігі
туындайды. Сондай-ақ академиялық адалдықты сақтау, бағалау критерийлерін
айқын ету және практикалық дағдыларды дамытуға бағытталған құралдарға
деген сұранысқа алып келеді {[}8{]}.

Осылайша, жаппай цифрлық платформалардың ыңғайлылығы мен «Деректер
базасы» пәнінің тақырыптары бойынша терең диагностика жүргізу
қажеттілігі арасында зерттеу мен әдістемелік алшақтық байқалады.
Қолданыстағы LMS жүйелерінде келесі мүмкіндіктер жеткіліксіз:

- жеке тақырыпшалар деңгейінде түсіндірілетін кері байланыс;

- студенттің профиліне тапсырманың күрделілігі мен форматын бейімдеу;

\begin{quote}
- тақырыптық «олқылықтарды» жедел анықтауға және курсқа мақсатты
\end{quote}

түзетулер енгізуге мүмкіндік беретін оқытушыларға арналған аналитикалық
жүйелер.

Қосымша шектеу ол сапалы тест базаларын қолмен әзірлеуге және типтік
қателіктерді талдауға арналған ресурстық шығындар болып табылады.

Осы міндеттерге жауап ретінде «Деректер базасы» пәні үшін жасанды
интеллект негізіндегі қателерді жіктеу моделі және жеке ұсыныстар
жүйесімен интеграцияланған интеллектуалды оқыту көмекшісі ұсынылады.
Дәстүрлі тесттік модульдерден айырмашылығы, жүйе тек жауаптарды бағалап
қана қоймайды, түсіндірмелі кері байланыс береді, жеке қайталау
траекторияларын қалыптастырады және әлсіз тақырыптарды (мысалы, қате
қосылулар, агрегаттау немесе фильтрациядағы қатені) анықтайды. Ал
оқытушы үшін тақырыптар мен топтар бойынша аналитикасы бар жинақталған
есептер қолжетімді болады. Бұл күнделікті талдауға кететін уақытты
үнемдеп, жетекшілікке арналған ресурстарды арттырады.

Ұсынылған тәсіл инженерлік іске асыруды (тестілеу модулі бар
веб‑қосымша) және зерттеу құрамдас бөліктерін (қателер классификаторын
оқыту және бағалау, тақырыптар бойынша дәлдік/толықтық көрсеткіштерін
талдау, түсіндірмелі кері байланыстың оқу нәтижелеріне әсерін зерттеу)
біріктіреді. Ғылыми тұрғыдан алғанда, маңыздысы - бағалау
процедураларының қайталануы, деректер мен көрсеткіштердің айқындығы,
сондай-ақ этикалық аспектілер - құпиялылық, ақпараттандырылған келісім,
модельдің ықтимал ауытқуларын аудиті {[}9{]}.

Практикалық жағынан, жүйе студенттердің ағымдағы бақылауға және
магистратураға түсу емтихандарына дайындық кезінде мүмкіндіктерін
арттыруға, сондай-ақ рутиналық тапсырмаларды автоматтандыру және
бақылаудың диагностикалық мүмкіндіктерін жақсарту арқылы оқытушылардың
жүктемесін азайтуға бағытталған.

{\bfseries Материалдар мен әдістер.} Жүйенің архитектурасы мен жобалануы.
Жүйе мобильді

бейімделуі бар веб-қосымша түрінде жүзеге асырылған. Фронтенд React
(TypeScript) тілінде жазылған, бэкенд -- Node.js, ML қызметтері - Python
(Flask) және PyTorch кітапханасын пайдалану арқылы жүзеге асырылған.
Логтар мен тест нәтижелерін сақтау үшін PostgreSQL деректер қорын
басқару жүйесі (ДҚБЖ) қолданылған, кэш үшін - Redis пайдаланылды.
Модельдерді оқыту NVIDIA Tesla K20 GPU-мен және 112 ГБ оперативті жады
бар серверде жүргізілген.

Тестілік модуль. Сұрақтар банкі SQL тақырыбы бойынша 250 тапсырма
енгізілген:

- DML операциялары (SELECT, INSERT, UPDATE);

- JOINS (INNER, LEFT, RIGHT, FULL);

- Агрегация (GROUP BY, HAVING, терезелік функциялар);

- теориялық сұрақтар (нормализация техникасы, индекстер).

Әрбір тапсырмаға дұрыс шешім үлгісі және типтік қателер сипаттамасы
берілген.

Қателер андатпасы және іріктеулерді дайындау.

Қолданушының енгізген SQL - сұраныстары мен ДҚБЖ қателер туралы
хабарламалары логтар арқылы автоматты түрде жиналды. Үш оқытушыдан
тұратын сараптамалық топ 5000 қателік мысалын келесі категориялар
бойынша қолмен бөліп белгіледі:

- синтаксистік қателер (кілттік сөздердің жетіспеушілігі, элементтердің
дұрыс

емес реті);

- семантикалық қателер (деректер кестесін/өрісін дұрыс таңдамау);

- логикалық (JOIN шарттарындағы немесе агрегаттық функциялардағы
қателер).

Деректер оқыту (80\%), валидация (10\%) және тестілеу (10\%) үлгілеріне
бөлінді.

Жауаптарды талдау моделі.

Қателерді жіктеу үшін BERT-based (RUBERT) моделі қолданылып, SQL-логтар
корпусында қайта оқытылды:

- кілт сөздер мен идентификаторларды ескере отырып, SQL сұрауларын
токенизациялау;

- BERT-ті fine-tuning бір толық байланысқан қабатпен (hidden\_size=768);

- AdamW оптимизаторы, learning rate=2e-5, batch\_size=16.

Сапаны бағалау: әрбір қателік категориясы бойынша precision, recall және
F1-score.

Ұсынымдық модуль.

Классификация нәтижелері мен студенттің үлгерім профилі негізінде, ең
күрделі тақырыптарды болжау үшін градиентті бустинг (XGBoost) модулі
жүзеге асырылды. Модельде келесі белгілер қолданылды:

- әр санат бойынша қателіктердің жиілігі;

- орташа жауап беру уақыты;

- ұқсас тапсырмаларды орындау тарихы.

\begin{quote}
Модель 80/10/10 белгілері бойынша оқытылды.
\end{quote}

Адаптивті тестілеу.

Адаптивтілік Item Response Theory (IRT) екі параметрлі логистикалық
моделіне негізделген алгоритммен қамтамасыз етілді:



Мұндағы
-
дискриминация параметрі,
-
сұрақтың күрделілік деңгейі,
-
студент қабілетінің көрсеткіші.

және

параметрлері максималды ықтималдық әдісімен оқыту деректер жиынында
бағаланды. Әр жауаптан кейін жүйе студенттің ағымдағы қабілет деңгейін

қайта есептеп, соған сәйкес күрделігіне оңтайлы келесі сұрақты таңдап
отырады. Алынған параметрлер арқылы IRT тәсілінің онлайн оқыту
жүйелерінде тиімді қолдануға болатындығын көрсетті, бұл Choi мен
McClenen еңбектерінде де атап өтілген {[}10{]}.

Тиімділікті бағалау.

Оқу үлгерімінің жақсаруы туралы гипотезаны тексеру үшін студенттер екі
топқа бөлінді - эксперименттік және бақылау топтары (әрқайсысы 60
студенттен).

- бақылау тобы: статикалық тесттілер банкін;

- эксперименттік топ: жасанды интеллект негізіндегі кері байланысы бар

бейімделетін тестілеу.

Орташа баллдардың айыршылығы Student t-критерийі (p<0,05) арқылы
анықталды. Сонымен қатар, студенттердің қанағаттану деңгейі SUS (System
Usability Scale) сауалнамасы арқылы жинақталды {[}11,12{]}.

Аппараттық-бағдарламалық платформа.

TestGen атты бақылау тесттерін автоматты түрде генерациялау жүйесі
таратылған аппараттық-бағдарламалық платформа негізінде құрылған. Ол
есептеу ресурстарын, деректерді сақтау жүйесін және бағдарламалық
қызметтер кешенін біріктіреді. Есептеу кластерінің өзегін Intel Xeon
процессорлары мен NVIDIA Tesla K20 графикалық үдеткіштері бар серверлер
құрайды. Бұл құрылғылар машиналық оқыту және табиғи тілді генерациялау
операцияларын жоғары өнімділікпен орындауға мүмкіндік береді.

Бағдарламалық бөлік микросервистік архитектура қағидаттарымен әзірленіп,
Kubernetes ортасында орналастырылған. Мұндай шешім жүйенің
масштабталуын, тұрақтылығын және қызметтердің оқшаулануын қамтамасыз
етеді. Docker контейнеризациясы әрбір компонентті жеке басқаруға
мүмкіндік береді: контентті импорттау және түрлендіру модулі, табиғи
тілді өңдеу қызметтері, LLM-воркерлері, деректердің қайталануын бақылау
және сапаны қамтамасыз ету жүйесі, тестілерді жинақтау және экспорттау
модулі. Компоненттердің өзара әрекеттесуі АРІ шлюзі және Redis хабарлама
жүйесі арқылы жүзеге асады, бұл жүктемені тиімді теңестіруге мүмкіндік
береді.

Осылайша, TestGen платформасы жоғары өнімді есептеу инфрақұрылымын,
заманауи контейнерді басқару құралдарын және модульдік микросервистік
архитектураны үйлестіре отырып, электронды оқыту жүйелеріне арналған
тест тапсырмаларын автоматты генерациялау міндетін тиімді шешуге
мүмкіндік береді.

\fig{i/image82}[]

{\bfseries 1-сурет. TestGen бақылау тестілерін автоматты түрде генерациялау
жүйесінің жалпы архитектурасы}

1-суретте оқыту платформаларына (мысалы, Moodle және Platonus)
интеграциялауға бағытталған TestGen автоматтандырылған бақылау
тестілерін генерациялау жүйесінің жалпыланған архитектурасы берілген.
Архитектура төрт деңгейден тұрады:

1. Пайдаланушы деңгейі - оқытушы жүйемен веб-интерфейс арқылы
әрекеттеседі. Бұл деңгейде аутентификация жүргізіледі, курс пен
дәрістер таңдалады, тест параметрлері бапталады, сондай-ақ
генерацияланған тапсырмаларды өңдеу және алдын ала қарау мүмкіндігі
беріледі.

2. Веб-интерфейс - пайдаланушы мен ішкі сервистер арасындағы байланысты
қамтамасыз етеді. Ол параметрлерді жинауға, тапсырмалардың орындалу
күйін көрсетуге және өңдеу құралдарын ұсынуға жауап береді.

3. Сервистік деңгей (Backend Services)- келесі негізгі құрамдас
бөліктерден тұрады:

- LMS Integration -- оқу үдерісін басқару жүйелерімен (Moodle/Platonus)
өзара

әрекеттесуді қамтамасыз ететін модуль, ол дәрістерді автоматты түрде
импорттауға және тесттерді экспорттауға мүмкіндік береді.

- Content Import -- әр түрлі форматтағы (PDF, DOCX, PPT, HTML) оқу

материалдарын жүктеуге және түрлендіруге мүмкіндік береді.

- Segmentation and Generation -- материалдарды мағынасына сай блоктарға
бөледі

және олардың негізінде жасанды интеллект модельдерінің көмегімен тест
сұрақтарын генерациялайды.

\begin{quote}
- Assembly and Export -- ең тиімді сұрақтарды таңдайды, тапсырмалар
банкін
\end{quote}

қалыптастырады және оны Moodle XML форматында немесе тікелей LMS
жүйесіне экспорттайды.

4. Сыртқы жүйелер мен деректер қоймалары:

\begin{quote}
- Data Store (PostgreSQL, Object Storage, Redis) -- бастапқы файлдарды,
мәтін
\end{quote}

сегменттерін, эмбеддингтерді, генерацияланған сұрақтар мен
метадеректерді сақтау үшін қолданылады.

\begin{quote}
- AI ішкі жүйесі - сұрақтарды генерациялау және олардың сапасын талдау
\end{quote}

модельдерінің жұмысын қамтамасыз етеді.

\begin{quote}
- LMS (Moodle, Platonus) -- дайын тесттерді оқу курстарына орналастыруға
арналған
\end{quote}

соңғы орта болып табылады.

Сызбадағы деректер ағындары оқу материалдарын жүктеуден бастап, оларды
өндеу, сұрақтарды генерациялау және дайын тесттерді LMS жүйесіне
жариялау кезеңіне дейінгі ақпарат қозғалысын бейнелейді.

Ұсынылған жүйенің практикалық қолданылуын көрсету үшін эксперименттік
топтағы нақты жағдайдың мысалын келтіреміз. Студентке агрегаттық
функциялар мен кестелерді біріктіруді қолдануға арналған тапсырма
берілді.2-суретте келесі SQL-сұранысы көрсетілген:

\fig{i/image83}[]

{\bfseries 2-сурет. Студенттің SQL сұрауындағы қатені анықтау мысалы}

Жүйе қате түрін автоматты түрде анықтайды және тақырыптық ұсыныс
жасайды. Бұл жағдайда қате өріс атауына (student\_id орнына students\_id
қате өріс атауы) байланысты семантикалық қате анықталды.

{\bfseries Нәтижелер және талқылау.} Қателерді классификациялау сапасы.
RUBERT моделі үш негізгі қате санатын жіктеу барысында жоғары дәлдік
көрсеткіштерін көрсетті. F1-score ең үлкен көрсеткіші логикалық
қателерді тану болып табылады, бұл JOIN және агрегаттық функциялардың
дұрыс емес шарттарын дұрыс анықтау үшін өте маңызды (1-кесте).

{\bfseries 1 - кесте. RuBERT моделі бойынша қателіктерді жіктеу сапасының
көрсеткіштері}

%% \begin{longtable}[]{@{}
%%   >{\raggedright\arraybackslash}p{(\linewidth - 6\tabcolsep) * \real{0.4117}}
%%   >{\centering\arraybackslash}p{(\linewidth - 6\tabcolsep) * \real{0.1825}}
%%   >{\centering\arraybackslash}p{(\linewidth - 6\tabcolsep) * \real{0.2029}}
%%   >{\centering\arraybackslash}p{(\linewidth - 6\tabcolsep) * \real{0.2029}}@{}}
%% \toprule\noalign{}
%% \begin{minipage}[b]{\linewidth}\centering
%% {\bfseries Қателер категориясы}
%% \end{minipage} & \begin{minipage}[b]{\linewidth}\centering
%% {\bfseries Precision}
%% \end{minipage} & \begin{minipage}[b]{\linewidth}\centering
%% {\bfseries Recall}
%% \end{minipage} & \begin{minipage}[b]{\linewidth}\centering
%% {\bfseries F1-score}
%% \end{minipage} \\
%% \midrule\noalign{}
%% \endhead
%% \bottomrule\noalign{}
%% \endlastfoot
%% Синтаксистік & 0,91 & 0,88 & 0,89 \\
%% Семантикалық & 0,87 & 0,85 & 0,86 \\
%% Логикалық & 0,89 & 0,90 & 0,90 \\
%% \end{longtable}

1-кестеде үш негізгі қате санатын жіктеу кезінде ruBERT моделі бойынша
жоғары көрсеткіштер ұсынылған.

Қате санаттары бойынша F1-score мәндері 3-суретте график түрінде де
көрсетілген, бұл қате түрлері арасындағы айырмашылықтарды көрнекі түрде
бағалауға мүмкіндік береді.

\fig{i/image84}[]

{\bfseries 3-сурет. Қате санаттары бойынша F1-score мәндері}

Ұсынымдық модулінің тиімділігі

XGBoost негізіндегі ұсыным жүйесі қайталап оқытуды қажет ететін ең
өзекті үш негізгі тақырыпты жоғары дәлдікпен болжайды.

Модельдің ROC-AUC және дәлдік (accuracy) көрсеткіштері төменде 2 кестеде
көрсетілген.

{\bfseries 2 - кесте. XGBoost негізіндегі ұсыныс модулінің сапа
көрсеткіштері}

%% \begin{longtable}[]{@{}
%%   >{\raggedright\arraybackslash}p{(\linewidth - 2\tabcolsep) * \real{0.5383}}
%%   >{\raggedright\arraybackslash}p{(\linewidth - 2\tabcolsep) * \real{0.4617}}@{}}
%% \toprule\noalign{}
%% \begin{minipage}[b]{\linewidth}\raggedleft
%% {\bfseries Көрсеткіштер}
%% \end{minipage} & \begin{minipage}[b]{\linewidth}\centering
%% {\bfseries Мәні}
%% \end{minipage} \\
%% \midrule\noalign{}
%% \endhead
%% \bottomrule\noalign{}
%% \endlastfoot
%% Accuracy & 0,84 \\
%% ROC-AUC & 0,91 \\
%% \end{longtable}

Ұсыным модулінің тиімділік көрсеткіштері 4-суретте толығырақ
көрсетілген.

\fig{i/image85}[]

{\bfseries 4-сурет. Ұсыным модулінің Accuracy және ROC-AUC көрсеткіштері}

Қанағаттану деңгейін бағалау (SUS)

Курсты аяқтағаннан кейін қатысушылар екі топқа бөлінді.3-кестеде
көрсетілгендей, ЖИ-негізіндегі кері байланысы бар адаптивті тестілеуді
пайдаланған эксперименттік топ бақылау тобына қарағанда айтарлықтай
жақсы нәтиже көрсетті.

{\bfseries 3 - кесте. Бақылау және эксперименттік топтардың көрсеткіштері}

%% \begin{longtable}[]{@{}
%%   >{\centering\arraybackslash}p{(\linewidth - 6\tabcolsep) * \real{0.2218}}
%%   >{\centering\arraybackslash}p{(\linewidth - 6\tabcolsep) * \real{0.3030}}
%%   >{\centering\arraybackslash}p{(\linewidth - 6\tabcolsep) * \real{0.3335}}
%%   >{\centering\arraybackslash}p{(\linewidth - 6\tabcolsep) * \real{0.1417}}@{}}
%% \toprule\noalign{}
%% \begin{minipage}[b]{\linewidth}\centering
%% {\bfseries Топ}
%% \end{minipage} & \begin{minipage}[b]{\linewidth}\centering
%% {\bfseries Орташа балл дейін (\%)}
%% \end{minipage} & \begin{minipage}[b]{\linewidth}\centering
%% {\bfseries Орташа балл кейін (\%)}
%% \end{minipage} & \begin{minipage}[b]{\linewidth}\centering
%% {\bfseries Жақсару}
%% \end{minipage} \\
%% \midrule\noalign{}
%% \endhead
%% \bottomrule\noalign{}
%% \endlastfoot
%% Бақылау & 64,3 & 68,1 & 3,8 \\
%% Эксперименттік & 63,8 & 81,6 & 17,8 \\
%% \end{longtable}

System Usability Scale шкаласы бойынша қатысушылар жүйенің ыңғайлылығы
мен пайдалылығын бағалады. Эксперименттік топ пайдаланушылық тәжірибенің
айтарлықтай артқанын атап өтті (4-кесте).

{\bfseries 4 - кесте. SUS шкаласы бойынша қанағаттану деңгейі}

%% \begin{longtable}[]{@{}
%%   >{\centering\arraybackslash}p{(\linewidth - 4\tabcolsep) * \real{0.3510}}
%%   >{\centering\arraybackslash}p{(\linewidth - 4\tabcolsep) * \real{0.3491}}
%%   >{\centering\arraybackslash}p{(\linewidth - 4\tabcolsep) * \real{0.2999}}@{}}
%% \toprule\noalign{}
%% \begin{minipage}[b]{\linewidth}\centering
%% {\bfseries Көрсеткіш}
%% \end{minipage} & \begin{minipage}[b]{\linewidth}\centering
%% {\bfseries Бақылау тобы}
%% \end{minipage} & \begin{minipage}[b]{\linewidth}\centering
%% {\bfseries Эксперименттік топ}
%% \end{minipage} \\
%% \midrule\noalign{}
%% \endhead
%% \bottomrule\noalign{}
%% \endlastfoot
%% SUS-score (0 --100) & 67,2 & 82,5 \\
%% \end{longtable}

Эксперименттік және бақылау топтарының SUS орташа мәндерінің
айырмашылығы 15,3 балды құрады және бұл айырмашылық статистикалық
тұрғыдан маңызды (p <{} 0,05). Бақылау және эксперименттік
топтардағы пайдаланушылық тәжірибені салыстыру нәтижелер 5-суретте
көрсетілген.

\fig{i/image86}[]

{\bfseries 5-сурет. Бақылау және эксперименттік топтардағы SUS көрсеткішін
салыстыру}

Стандартты LMS жүйелерінен айырмашылығы (мысалы, Moodle), онда тестілеу
көбінесе статикалық сипатта болып, тек бинарлық бағалаумен (дұрыс/бұрыс)
шектелетін болса, ұсынылған жүйе келесі мүмкіндіктерді қамтамасыз етеді:

- қате түрлерін автоматты түрде жіктеу;

- жекелендірілген ұсыныстар беру;

- күрделілік деңгейін адаптивті реттеу (IRT);

- ішкі тақырыптар бойынша аналитикалық панель.

5-кестеде функционалдық сипаттамаларды салыстыру келтірілген.

{\bfseries 5 - кесте. Ұсынылған жүйені дәстүрлі LMS-пен салыстыру}

%% \begin{longtable}[]{@{}
%%   >{\raggedright\arraybackslash}p{(\linewidth - 4\tabcolsep) * \real{0.4480}}
%%   >{\centering\arraybackslash}p{(\linewidth - 4\tabcolsep) * \real{0.2932}}
%%   >{\centering\arraybackslash}p{(\linewidth - 4\tabcolsep) * \real{0.2589}}@{}}
%% \toprule\noalign{}
%% \begin{minipage}[b]{\linewidth}\centering
%% {\bfseries Сипаттама}
%% \end{minipage} & \begin{minipage}[b]{\linewidth}\centering
%% {\bfseries Дәстүрлі LMS}
%% \end{minipage} & \begin{minipage}[b]{\linewidth}\centering
%% {\bfseries Ұсынылған жүйе}
%% \end{minipage} \\
%% \midrule\noalign{}
%% \endhead
%% \bottomrule\noalign{}
%% \endlastfoot
%% Қате түрлерін диагностикалау & Жоқ & Иә (ruBERT) \\
%% Жекелендірілген ұсыныстар & Шектеулі & Иә (XGBoost) \\
%% Күрделіліктің адаптивтілігі & Жоқ & Иә (IRT) \\
%% Ішкі тақырыптық аналитика & Ішінара & Толық \\
%% Түсіндірмесі бар кері байланыс & Жоқ & Иә \\
%% \end{longtable}

Алынған деректер жасанды интеллект модельдерін жауаптарды талдау және
бейімделген тестілеу жүйесіне енгізу «Деректер базасы» пәнін оқыту
тиімділігін едәуір арттыратынын көрсетеді {[}13{]}. Қателерді жіктеу
моделінің жоғары көрсеткіштері студенттердің проблемалық тұстарын дәл
диагностикалауға мүмкіндік береді, ал ROC-AUC = 0.91 мәні бар ұсынымдық
модулі қайта қарауды қажет ететін тақырыптарды сенімді түрде анықтайды.
Сонымен қатар, ұсынылған жүйе жекелендіру деңгейі, диагностиканың
тереңдігі және адаптивтілігі бойынша дәстүрлі шешімдерден айтарлықтай
басым түседі.

Университет ортасында нақты пайдалану кезінде диагностика мен
жекелендіру сапасының өсуін көрсететін үш интеграцияланған құрамдас
бөлігі бар SQL кешенді ai оқыту жүйесі ұсынылған. Қателердің түрлерін
түсіндірмелі бөлу, тақырыптар бойынша жеке ұсыныстар, IRT бойынша
күрделілікті реттеу сияқты компоненттер іс жүзінде маңызды. Әдістеме мен
архитектуралық шешімдер ат циклінің сабақтас пәндеріне көшіріледі және
масштабталады. Әрі қарайғы зерттеуде датасет пен курстарды кеңейту,
белгілеудің жартылай/нашар бақыланатын әдістерін енгізу,
эмоционалды-танымдық факторлардың әсерін зерттеу және прогресстің
жетілдірілген аналитикасы жоспарлануда.

Ұсынылған жүйе «Деректер базасы» пәнін және соған іргелес IT-бағыттарды
оқытуда маңызды педагогикалық мәнге ие. Біріншіден, қателерді
автоматтандырылған түрде жіктеу білімді формалды тексеруден
диагностикалық бағалауға көшуге мүмкіндік береді. Бұл тәсіл ішкі
тақырыптар (JOIN, агрегаттау, ішкі сұраныстар және т.б.) бойынша типтік
қиындықтарды анықтауға бағытталған. Нәтижесінде бақылаудың
қалыптастырушы функциясы күшейіп, мазмұнды әрі түсіндірмелі кері
байланыс қамтамасыз етіледі.

Екіншіден, ұсынымдық модульдің интеграциясы оқу траекториясын
жекелендіруге ықпал етеді: студент нақты тақырыптарды қайталауға
бағытталған мақсатты ұсыныстар алады. Бұл оқытудың саналылығын арттырып,
дайындықтың үзік-үзік болуын азайтады.

Үшіншіден, IRT-адаптивті тестілеуді қолдану білім алушының ағымдағы
деңгейіне байланысты тапсырмалардың күрделілігін реттеуге мүмкіндік
береді. Соның нәтижесінде білім алушының мүмкіндігіне сай тиімді оқу
аймағы сақталып, тапсырмалардың шамадан тыс қиын немесе тым жеңіл
болуынан туындайтын демотивацияны азайтады.

Зерттеудің әдістемелік үлесі -- жасанды интеллект құралдарын қолданбалы
сипаттағы пәнге интеграциялаудың қайта қолдануға болатын моделін
әзірлеуінде. Мұнда тек теориялық білімді тексеру ғана емес, сонымен
қатар практикалық шешім құрастыру дағдыларын диагностикалау маңызды.
Ұсынылған тәсіл АТ-циклдің басқа пәндеріне масштабтауға, сондай-ақ
қорытынды аттестаттау мен қабылдау емтихандарына дайындық барысында
қолдануға мүмкіндік береді.

{\bfseries Қорытынды.} Деректер базасын меңгеру сапасы түлектердің
қолданбалы міндеттерді шешу қабілетіне тікелей әсер етеді, ал
университеттер үшін еңбек нарығының жылдам өзгеріп жатқан жағдайында
білім беру бағдарламаларының өзектілігін сақтаудың негізгі факторы болып
табылады. Сондықтан заманауи интеллектуалды технологияларды
университеттік пәндерге енгізу жай ғана үрдіске ілесу емес, оқытудың
сапасы мен тұрақтылығын арттырудың пәрменді құралы болып табылады
мысалы: жасанды интеллект модельдері арқылы тақырыптар мен тақырыпшалар
бойынша нақты диагностика жасауға, траекторияны персонализациялауға
(тапсырмалардың күрделілігімен форматын бейімдеу), түсіндірмелі кері
байланыс ұсыну және оқытушыға курсқа жедел түзетулер енгізу үшін
аналитикалық панельдер беру. Нәтижесінде оқу бағдарламалары тартымды әрі
бәсекеге қабілетті бола түседі, ал студенттер қысқа мерзімді курстарға
тән «жылдам» дағдылардан бөлек, ұзақ мерзімді кәсіби табысқа негіз
болатын берік және өзекті білім алады.

Бұл мақалада SQL оқытуға арналған ЖИ-негізіндегі жүйе сипатталып, оның
тиімділігі эксперименттік түрде бағаланды. Жүйе RUBERT-қате
классификациясын, XGBoost-ұсынымдық модулін және IRT-адаптивті
тестілеуін біріктіреді.120 студент деректеріне сүйене отырып, жүйе
келесі нәтижелер көрсетті:

- қате түрлері бойынша жоғары диагностикалық дәлдік (F1 0,90-ға дейін);

- тақырыптарды қайталау үшін сенімді ұсынымдар (ROC-AUC 0,91, accuracy

0,84);

- пайдаланушылар тәжірибесінің айтарлықтай жоғары деңгейі (SUS 82.5
қарсы

67.2).

Бұл нәтижелер «диагностика-ұсыным-адаптивтілік» байланысының SQL-ді
жекелендірілген оқытудағы практикалық, тиімділігін дәлелдейді.

Зерттеудің шектеулеріне бір оқу бағдарламасының студенттерін ғана қамту,
қателерді қолмен белгілеудің еңбек сыйымдылығы, сондай-ақ тек SQL-ге
шоғырлану жатады. Болашақта жоспарланған қадамдар: зерттеуді кеңейту,
белгілерді арзандату үшін жартылай немесе әлсіз бақыланатын әдістерді
енгізу, сондай-ақ эмоция-танымдық факторлардың тест нәтижелеріне әсерін
зерттеу. Бұл жұмыс негізінен SQL оқытуға арналған жасанды интеллект
жүйесінің архитектурасы мен инженерлік-пайдалану аспектілеріне
шоғырланады: қабаттар мен интерфейстер, сенімділік пен масштабталуды
қамтамасыз ету.

{\bfseries Әдебиеттер}

1. Gaitantzi A., Kazanidis I. The role of artificial intelligence in
computer science education: A systematic review with a focus on
database instruction// Applied Sciences.-2025.-Vol.15(7):3960 DOI
10.3390/app15073960.

2. Ocen S., Elasu J., Aarakit S.M., Olupot C. Artificial intelligence in
higher education institutions: review of innovations, opportunities
and challenges // Frontiers in Education.-2025. -Vol.10.-P.1-34.
DOI~\href{https://doi.org/10.3389/feduc.2025.1530247}{10.3389/feduc.2025.1530247}.

3. Fabijanić М., Mekterović I. Partial SQL query assessment. Conference:
46th MIPRO ICT and Electronics Convention (MIPRO), 2023.
DOI~\href{https://doi.org/10.23919/MIPRO57284.2023.10159706}{10.23919/MIPRO57284.2023.10159706}.

4. Schwartz D.R., Rivas P. An Automated SQL Query Grading System Using An
Attention-Based Convolutional Neural Network. arXiv preprint, 2024.
DOI 10.48550/arXiv.2406.15936.

5. Maher Joe Khan Omar Jian. Personalized learning through AI // Advances
in Engineering Innovation. -2023. Vol.5(1). -P.16-19.
DOI~\href{https://doi.org/10.54254/2977-3903/5/2023039}{10.54254/2977-3903/5/2023039}.

6. Jose B., Cherian J., Verghis A., Varghise S., Mumthas S., Sibichan J.
The cognitive paradox of AI in education: between enhancement and
erosion//Front Psychol.-2025.-Vol.16:1550621.
\href{https://doi.org/10.3389/fpsyg.2025.1550621}{DOI
10.3389/fpsyg.2025.1550621}.

7. Zaman B.U. Transforming Education Through AI Benefits Risks and
Ethical Considerations //TechRxiv.-2023.-P.1-16. DOI
\href{https://doi.org/10.36227/techrxiv.24231583.v1}{10.36227/techrxiv.24231583.v1}.

8. Бермус А., Сериков В., Алтыникова Н. Содержание педагогического
образования в современном мире: смыслы, проблемы, практики и
перспективы развития // Вестник РУДН. Серия: Психология и педагогика.
-2021.-Vol.18(4).-C.667-691. DOI 10.22363/2313-1683-2021-18-4-667-691.

9. Мантуленко В. Влияние искусственного интеллекта на успеваемость,
познавательную активность и качество обучения студентов //
Научно-методический электронный журнал «Концепт».-2025.- № 06.-
С.251-272. DOI 10.24412/2304-120X-2025-11117.

10. Choi Y., McClenen C. Development of Adaptive Formative Assessment
System Using Computerized Adaptive Testing and Dynamic Bayesian Networks
//~Applied Sciences. - 2020. - Vol.10(22): 8196. DOI
\href{https://doi.org/10.3390/app10228196}{10.3390/app10228196}.

11. Dabingaya M. Analyzing the Effectiveness of AI-Powered Adaptive
Learning Platforms in Mathematics Education //~Interdisciplinary Journal
Papier Human Review.-2022.-Vol.3(1).- P.1-7. DOI
\href{https://doi.org/10.47667/ijphr.v3i1.226}{10.47667/ijphr.v3i1.226}.

12. Тагирова Л.Ф., Зубкова Т.М. Интеллектуальная система адаптивного
тестирования // Научно-технический вестник информационных технологий,
механики и оптики. -2023.- Т.23(4). - C.757-766. DOI
10.17586/2226-1494-2023-23-4-757-766.

13. Zhou J., Zhang J., Li H. Exploring the Use of Artificial
Intelligence in Teaching Management and Evaluation Based on Citation
Space Analysis // Journal of Education and Educational
Research.-2023.-Vol.3(2).- P.42-45.
DOI~\href{https://doi.org/10.54097/jeer.v3i2.9014}{10.54097/jeer.v3i2.9014}.

{\bfseries References}

1. Gaitantzi A., Kazanidis I. The role of artificial intelligence in
computer science education: A systematic review with a focus on database
instruction// Applied Sciences.-2025.-Vol.15(7). -P.1-34. DOI
10.3390/app15073960.

2. Ocen S., Elasu J., Aarakit S.M., Olupot C. Artificial intelligence in
higher education institutions: review of innovations, opportunities and
challenges // Frontiers in Education.-2025. -Vol.10.-P.1-34.
DOI~\href{https://doi.org/10.3389/feduc.2025.1530247}{10.3389/feduc.2025.1530247}.

3. Fabijanić М., Mekterović I. Partial SQL query assessment. Conference:
46th MIPRO ICT and Electronics Convention (MIPRO), 2023.
DOI~\href{https://doi.org/10.23919/MIPRO57284.2023.10159706}{10.23919/MIPRO57284.2023.10159706}.

4. Schwartz D.R., Rivas P. An Automated SQL Query Grading System Using
An Attention-Based Convolutional Neural Network. arXiv preprint, 2024.
DOI 10.48550/arXiv.2406.15936.

5. Maher Joe Khan Omar Jian. Personalized learning through AI //
Advances in Engineering Innovation. -2023. Vol.5(1). -P.16-19.
DOI~\href{https://doi.org/10.54254/2977-3903/5/2023039}{10.54254/2977-3903/5/2023039}.

6. Jose B., Cherian J., Verghis A., Varghise S., Mumthas S., Sibichan J.
The cognitive paradox of AI in education: between enhancement and
erosion//Front Psychol.-2025.-Vol.16:1550621.
\href{https://doi.org/10.3389/fpsyg.2025.1550621}{DOI
10.3389/fpsyg.2025.1550621}.

7. Zaman B.U. Transforming Education Through AI Benefits Risks and
Ethical Considerations //TechRxiv.-2023.-P.1-16. DOI
\href{https://doi.org/10.36227/techrxiv.24231583.v1}{10.36227/techrxiv.24231583.v1}.

8. Bermus A., Serikov V., Altynikova N. Soderzhanie pedagogicheskogo
obrazovanija v sovremennom mire: smysly, problemy, praktiki i
perspektivy razvitija // Vestnik RUDN. Serija: Psihologija i pedagogika.
-2021.-Vol.18(4).-C.667-691. DOI
10.22363/2313-1683-2021-18-4-667-691.{[}in Russian{]}

9. Mantulenko V. Vlijanie iskusstvennogo intellekta na
uspevaemost', poznavatel' nuju
aktivnost'{}

i kachestvo obuchenija studentov // Nauchno-metodicheskij jelektronnyj
zhurnal «Koncept».

2025. - № 06.- S.251-272. DOI 10.24412/2304-120X-2025-11117. {[}in
Russian{]}

10. Choi Y., McClenen C. Development of Adaptive Formative Assessment
System Using Computerized Adaptive Testing and Dynamic Bayesian Networks
//~Applied Sciences. - 2020. - Vol.10(22): 8196. DOI
\href{https://doi.org/10.3390/app10228196}{10.3390/app10228196}.

11. Dabingaya M. Analyzing the Effectiveness of AI-Powered Adaptive
Learning Platforms in Mathematics Education //~Interdisciplinary Journal
Papier Human Review.-2022.-Vol.3(1).- P.1-7. DOI
\href{https://doi.org/10.47667/ijphr.v3i1.226}{10.47667/ijphr.v3i1.226}.

12. Tagirova L.F., Zubkova T.M. Intellektual' naja
sistema adaptivnogo testirovanija // Nauchno-tehnicheskij vestnik
informacionnyh tehnologij, mehaniki i optiki. -2023.- T.23(4). -
C.757-766. DOI 10.17586/2226-1494-2023-23-4-757-766. {[}in Russian{]}

13. Zhou J., Zhang J., Li H. Exploring the Use of Artificial
Intelligence in Teaching Management and Evaluation Based on Citation
Space Analysis // Journal of Education and Educational
Research.-2023.-Vol.3(2).- P.42-45.
DOI~\href{https://doi.org/10.54097/jeer.v3i2.9014}{10.54097/jeer.v3i2.9014}.

\emph{{\bfseries Авторлар туралы мәлімет}}

Жарлықасов Б.Ж. - жаратылыстану ғылымдарының магистрі, Ахмет
Байтұрсынұлы атындағы Қостанай өңірлік университеті, Қостанай,
Қазақстан, е-mail:
bakhtiyarzbj@gmail.com;

Баенова Г.М. - PhD, Л.Н. Гумилев атындағы Еуразия ұлттық университеті,
Астана, Қазақстан
е-mail:gulmmira@mail.ru;

Мауленов Қ.С. - PhD, Ахмет Байтұрсынұлы атындағы Қостанай өңірлік
университеті, Қостанай, Қазақстан, е-mail:
k.maulenov070693@gmail.com;

Сыздыкова А.М. - жаратылыстану ғылымдарының магистрі, Л.Н. Гумилев
атындағы Еуразия ұлттық университеті, Астана, Қазақстан,
е-mail:aierke91@mail.ru.

\emph{{\bfseries Information about the authors}}

Zharlykassov B.Zh. - Master of Natural Sciences, Akhmet Baitursynuly
Kostanay Regional University, Kostanay, Kazakhstan, е-mail:
bakhtiyarzbj@gmail.com;

Baenova G.M. - PhD, L.N. Gumilyov Eurasian National University, Astana,
Kazakhstan, е-mail:
gulmmira@mail.ru;

Maulenov K.S. - PhD, Akhmet Baitursynuly Kostanay Regional University,
Kostanay, Kazakhstan, е-mail:
k.maulenov070693@gmail.com;

Syzdykova A.M. - Master of Natural Sciences, L.N. Gumilyov Eurasian
National University, Astana, Kazakhstan, е-mail:
aierke91@mail.ru.\